% Le linee di attività di Halue come colonne a gradini, in ordine di peso, sulla base
% comune della consulenza tecnologica. Il disegno è largo quanto il testo (150 mm).
\begin{tikzpicture}[every node/.style={text=FiguraTesto}]
	\path[intestazione] (0, 0) rectangle (150mm, 9mm);
	\node[titolo] at (75mm, 4.5mm) {Consulenza tecnologica};
	\foreach \i/\altezza/\titolo/\dettaglio in {
		1/33/{1. \emph{E-commerce}}/{Shopify e Salesforce Commerce Cloud, B2C e B2B},
		2/30/{2. Integrazione di sistemi}/{ERP, logistica, servizi di terze parti, iPaaS},
		3/27/{3. Intelligenza artificiale}/{dalla valutazione della prontezza agli agenti},
		4/24/{4. CRM}/{Salesforce},
		5/21/{5. Formazione}/{metodi, piattaforme e tecnologia},
		6/18/{6. Interfaccia ed esperienza utente}/{}} {
		\pgfmathsetlengthmacro\x{(\i-1)*(22.5mm+3mm)}
		\ifnum\i=3
			\path[draw=FiguraAccento, line width=0.8pt, fill=white, rounded corners=3pt]
				(\x, 11mm) rectangle ++(22.5mm, \altezza mm);
		\else
			\path[draw=FiguraTenue, fill=white, rounded corners=3pt]
				(\x, 11mm) rectangle ++(22.5mm, \altezza mm);
		\fi
		\node[anchor=north, text width=21mm, align=center, inner sep=0pt,
			font=\fontsize{7.5}{9}\selectfont,
			execute at begin node={\hyphenpenalty=10000\exhyphenpenalty=10000}]
			at ($(\x, 11mm) + (11.25mm, \altezza mm - 2mm)$) {\voce{\titolo}{\dettaglio}};
	}
	% Una parentesi bordeaux mostra dove entra l'intelligenza artificiale.
	\draw[FiguraAccento, line width=0.6pt] (0, 46mm) -- ++(0, 1mm) -- (48mm, 47mm) -- ++(0, -1mm);
	\node[anchor=south, inner sep=0pt, font=\fontsize{7.5}{9}\selectfont, text=FiguraAccento]
		at (24mm, 48.2mm) {con l'intelligenza artificiale};
\end{tikzpicture}
